\documentclass[11pt]{amsart}

\usepackage[margin=1.15in]{geometry}
\usepackage{amsmath,amssymb,amsthm,mathtools}
\usepackage{enumitem}
\usepackage{microtype}
\usepackage{hyperref}
\usepackage[nameinlink,capitalize]{cleveref}

\hypersetup{
  colorlinks=true,
  linkcolor=blue,
  citecolor=blue,
  urlcolor=blue
}

\theoremstyle{plain}
\newtheorem{theorem}{Theorem}[section]
\newtheorem{proposition}[theorem]{Proposition}
\newtheorem{lemma}[theorem]{Lemma}
\newtheorem{corollary}[theorem]{Corollary}

\theoremstyle{definition}
\newtheorem{definition}[theorem]{Definition}
\newtheorem{hypothesis}[theorem]{Hypothesis}

\theoremstyle{remark}
\newtheorem{remark}[theorem]{Remark}

\numberwithin{equation}{section}

\newcommand{\R}{\mathbb R}
\newcommand{\N}{\mathbb N}
\newcommand{\eps}{\varepsilon}
\newcommand{\loc}{\mathrm{loc}}
\newcommand{\supp}{\operatorname{supp}}
\newcommand{\divergence}{\operatorname{div}}
\newcommand{\dist}{\operatorname{dist}}
\newcommand{\esssup}{\operatorname*{ess\,sup}}
\newcommand{\calS}{\mathcal S}
\newcommand{\calR}{\mathcal R}
\newcommand{\calX}{\mathcal X}
\newcommand{\calU}{\mathcal U}

\title[Type I blow-up limits and residual classes]
{Type I Blow-up Limits and a Residual-Class Criterion for the Three-Dimensional Navier--Stokes Equations}

\author{Guillem Duran Ballester}
\email{guillem@fragile.tech}
\date{}

\subjclass[2020]{35Q30, 35B44, 35B40, 76D05}
\keywords{Navier--Stokes equations, Type I singularities, ancient solutions, suitable weak solutions, Seregin limits, Liouville theorems}

\begin{document}

\begin{abstract}
We prove a detailed reduction from a Paper--0-admissible Type I singular
point of a finite-energy suitable weak solution of the three-dimensional
Navier--Stokes equations to a normalized nonzero bounded ancient solution of
the centered self-similar Navier--Stokes equation.  The extraction keeps track
of pressure gauges, local compactness, stability of the local energy
inequality, and the velocity concentration supplied by the local
Caffarelli--Kohn--Nirenberg analysis, yielding a smooth bounded ancient
profile for the centered equation on \(\R^3\times\R\), together with a
quantitative velocity nonvanishing normalization.

We then organize the possible normalized Seregin ancient limits into
small-amplitude, stationary \(L^3\), uniformly \(L^3\)-tight,
structured/decay, and residual classes.  The small-amplitude branch is
excluded by a bounded ancient perturbative Liouville theorem proved below, and
the stationary \(L^3\) branch is excluded by the theorem of
Ne\v{c}as--R\r{u}\v{z}i\v{c}ka--\v{S}ver\'ak.  Assuming Liouville criteria
for the tight and structured/decay classes, exclusion of Paper--0-admissible
finite-energy Type I sources is reduced to one remaining assertion: the
residual class of normalized Seregin limits generated by Paper--0-admissible
Type I sources is empty, or equivalently satisfies a replacement rigidity
theorem.
\end{abstract}

\maketitle

\section{Introduction and main results}\label{sec:introduction}

\subsection{Type I blow-up limits}

Let \(u:\R^3\times(0,T)\to\R^3\) solve the incompressible Navier--Stokes
equations
\begin{equation}\label{eq:NS}
   \partial_tu+(u\cdot\nabla)u+\nabla p=\Delta u,\qquad
   \divergence u=0.
\end{equation}
The local regularity theory of Caffarelli--Kohn--Nirenberg \cite{CKN1982}
detects singular points through scale-invariant concentration of velocity and
pressure.  A terminal point \(z_*=(x_*,T)\) is of Type I if
\begin{equation}\label{eq:type-I-bound}
   \limsup_{t\uparrow T}\sqrt{T-t}\,
   \|u(t)\|_{L^\infty(\R^3)}<\infty .
\end{equation}
The blow-up extraction in this paper uses a selected velocity concentration
sequence from the local concentration paper and assumes only that this
velocity concentration does not escape into the terminal slice after
rescaling.  This is the Paper--0-admissible source notion defined precisely
in \cref{def:source}.  It replaces the stronger all-scale compactness
assumption on scale-invariant energy, pressure, and dissipation
quantities by the exact non-escape condition needed for nontriviality of the
ancient limit.

\subsection{Normalized Seregin limits}

For scales \(\lambda\downarrow0\), the parabolic rescaling about \(z_*\) is
\[
   u^{(\lambda)}(x,t)
   =
   \lambda u(x_*+\lambda x,T+\lambda^2t),\qquad
   p^{(\lambda)}(x,t)
   =
   \lambda^2p(x_*+\lambda x,T+\lambda^2t).
\]
The rescaled time intervals exhaust \((-\infty,0)\).  The Type I bound
becomes a uniform bound
\[
   \|u^{(\lambda)}(t)\|_{L^\infty(\R^3)}
   \le M(-t)^{-1/2}
\]
on compact subintervals of \((-\infty,0)\), and compactness produces an
unrenormalized ancient solution \((U,Q)\).  The associated centered variables
are
\[
   y=\frac{x}{\sqrt{-t}},\qquad
   \tau=-\log(-t),\qquad
   V(y,\tau)=\sqrt{-t}\,U(y\sqrt{-t},t).
\]
They transform the ancient Navier--Stokes solution into the centered equation
\begin{equation}\label{eq:centered}
   \partial_\tau V+(V\cdot\nabla)V+\nabla\Pi
   =
   \Delta V-\frac12(V+y\cdot\nabla V),
   \qquad \divergence V=0.
\end{equation}
The selected velocity concentration at the singular point persists through
the extraction and gives a compact cylinder \(K\Subset\R^3\times\R\) and a
constant \(\eta_0>0\) such that
\begin{equation}\label{eq:nonvanishing}
   \iint_K |V|^3\,dy\,d\tau\ge\eta_0 .
\end{equation}

\subsection{Branch alternatives and residual obstruction}

No unconditional Liouville theorem is known for all bounded ancient solutions
of \eqref{eq:centered}.  The later part of the paper therefore records a
finite class decomposition needed for a conditional Type I exclusion
criterion.  A normalized Seregin limit is placed into one of the following
broad classes:
\[
   \calX_{\mathrm{small}},\qquad
   \calX_{\mathrm{stat}\text{-}L^3},\qquad
   \calX_{L^3\mathrm{-tight}},\qquad
   \calX_{\mathrm{sd}}(I,J),\qquad
   \calR(\calS;I,J).
\]
The first four are non-residual classes: small amplitude, stationary \(L^3\),
uniformly \(L^3\)-tight, and structured/decay.  The residual class is the
complement of their union in a fixed Seregin collection \(\calS\).  The
residual statement is therefore not an extra branch; it is the assertion that
no normalized Seregin limit can avoid all specified Liouville criteria.

\subsection{Main theorems}

\begin{theorem}[Normalized Type I blow-up limit]\label{thm:extraction}
Let \((u,p,z_*)\) be a Paper--0-admissible finite-energy Type I source.  Then there
exist scales \(\lambda_k\downarrow0\), pressure gauges \(a_k(t)\), an ancient
suitable weak solution \((U,Q)\) on \(\R^3\times(-\infty,0)\), and a centered
profile \((V,\Pi)\) on \(\R^3\times\R\), such that, after passing to a
subsequence,
\[
   u_k\to U
   \quad\hbox{strongly in }L^q_{\loc}(\R^3\times(-\infty,0)),
   \qquad 1\le q<\infty,
\]
and
\[
   p_k-a_k(t)\rightharpoonup Q
   \quad\hbox{weakly in }L^{3/2}_{\loc}(\R^3\times(-\infty,0)).
\]
The centered field \(V\) is smooth, bounded, solves \eqref{eq:centered}, and
satisfies the velocity nonvanishing normalization \eqref{eq:nonvanishing}.
\end{theorem}

The branch definitions give the exhaustive alternative
\[
   V\in
   \calX_{\mathrm{small}}
   \cup
   \calX_{\mathrm{stat}\text{-}L^3}
   \cup
   \calX_{L^3\mathrm{-tight}}
   \cup
   \calX_{\mathrm{sd}}(I,J)
   \cup
   \calR(\calS;I,J).
\]
This is recorded later as \cref{prop:exhaustion}.  It is a direct consequence
of the definition of the residual class, not an independent Liouville theorem.

\begin{theorem}[Residual-class criterion for Type I exclusion]
\label{thm:residual-criterion}
Fix admissible parameter sets \(I,J\).  Assume
\cref{hyp:tight-liouville,hyp:structured-decay-liouville,hyp:no-residual} for
normalized Seregin limits generated by Paper--0-admissible finite-energy Type
I sources.  Then no Paper--0-admissible finite-energy Type I source exists.
\end{theorem}

Thus the only remaining analytic hypotheses needed to upgrade the criterion to a
full Type I exclusion theorem are precisely the stated tight,
structured/decay, and residual Liouville assertions.

The dependency chain is
\[
\begin{array}{c}
\text{Paper--0-admissible finite-energy Type I source}
\\
\Downarrow
\\
\text{nonzero normalized Seregin limit}
\\
\Downarrow
\\
\text{small/stat/tight/structured/residual alternative}
\\
\Downarrow
\\
\text{contradiction under tight, structured/decay, and no-residual hypotheses.}
\end{array}
\]

\[
\begin{array}{c|c|c}
\text{Branch} & \text{Status in Paper I} & \text{Reference or source}\\ \hline
\text{small amplitude} & \text{proved} & \text{Section~\ref{sec:liouville}}\\
\text{stationary }L^3 & \text{proved} & \text{\cite{NRS1996}}\\
\text{uniformly }L^3\text{-tight} & \text{hypothesis} &
   \text{Hypothesis~\ref{hyp:tight-liouville}}\\
\text{structured/decay} & \text{hypothesis} &
   \text{Hypothesis~\ref{hyp:structured-decay-liouville}}\\
\text{residual} & \text{hypothesis} & \text{residual rigidity problem}
\end{array}
\]

\subsection{Dependencies and conditional inputs}

The proof of the extraction theorem uses the velocity concentration
corollary from the local concentration paper, stated here as
\cref{prop:paper0-velocity-concentration}; this is the only prior result from
the series used in the blow-up extraction.  The tight-class,
structured/decay, and residual assertions are not used to construct the
ancient limit.  They enter only in the final conditional criterion as the
explicit hypotheses
\cref{hyp:tight-liouville,hyp:structured-decay-liouville,hyp:no-residual}.
The present paper contains the normalization, the class decomposition, and the
residual-class criterion; it does not claim an unconditional Type I exclusion
theorem without those Liouville hypotheses.

\subsection{Organization of the paper}

\Cref{sec:prelim} records suitable solutions, scale-invariant quantities,
source admissibility, scaling, pressure gauges, and local concentration.
\Cref{sec:blowup-compactness} proves the compactness theorem for Type I
rescalings away from the terminal time and constructs the ancient suitable
weak limit.  \Cref{sec:nontriviality} proves persistence of concentration and
nontriviality.  \Cref{sec:seregin} derives the centered equation and defines
normalized Seregin limits and collections.  \Cref{sec:classes} defines the
branch classes and the residual class.  \Cref{sec:liouville} records the
Liouville hypotheses for non-residual branches.  \Cref{sec:criterion} proves the
residual-class criterion.  Appendix \ref{app:no-cascade} records the absence
of additional centered dilation and translation symmetries.

\section{Suitable solutions, scale-invariant quantities, and local concentration}
\label{sec:prelim}

For \(z_0=(x_0,t_0)\), write
\[
   Q_r(z_0)=B_r(x_0)\times(t_0-r^2,t_0).
\]
When \(z_0=(0,0)\), write \(Q_r=Q_r(0,0)\).  The spatial average of \(f\) over
\(B_r(x_0)\) is
\[
   (f)_{B_r(x_0)}(t)=\frac1{|B_r|}\int_{B_r(x_0)}f(x,t)\,dx.
\]
Throughout the paper, \(\mathbb P\) denotes the Helmholtz--Leray projection
onto divergence-free vector fields on \(\R^3\).

\begin{definition}[Terminal singular point]\label{def:terminal-singular}
Let \((u,p)\) be a suitable weak solution on \(\R^3\times(0,T)\).  A terminal
point \(z_*=(x_*,T)\) is regular if there are \(r>0\) and \(K<\infty\) such
that
\[
   \esssup_{Q_r(z_*)}|u|\le K.
\]
It is singular otherwise.
\end{definition}

\begin{definition}[Suitable weak solution]\label{def:suitable}
Let \(\Omega\subset\R^3\times\R\) be open.  A pair \((u,p)\) is a suitable
weak solution of \eqref{eq:NS} in \(\Omega\) if
\[
   u\in L^\infty_tL^2_x(\Omega')\cap L^2_t\dot H^1_x(\Omega'),
   \qquad
   p\in L^{3/2}(\Omega')
\]
for every compact cylinder \(\Omega'\Subset\Omega\), if \((u,p)\) solves
\eqref{eq:NS} distributionally, \(\divergence u=0\), and if the local energy
inequality
\begin{align}\label{eq:lei}
   &\int_{\R^3}|u(x,t_2)|^2\phi(x,t_2)\,dx
   +2\int_{t_1}^{t_2}\int_{\R^3}|\nabla u|^2\phi\,dx\,dt  \notag\\
   &\le
   \int_{\R^3}|u(x,t_1)|^2\phi(x,t_1)\,dx
   +\int_{t_1}^{t_2}\int_{\R^3}|u|^2(\partial_t\phi+\Delta\phi)\,dx\,dt
   \notag\\
   &\quad
   +\int_{t_1}^{t_2}\int_{\R^3}(|u|^2+2p)u\cdot\nabla\phi\,dx\,dt
\end{align}
holds for every non-negative \(\phi\in C_c^\infty(\Omega)\) and almost every
\(t_1<t_2\) for which the displayed terms are defined.
\end{definition}

This is the Caffarelli--Kohn--Nirenberg notion of suitability
\cite{CKN1982}; the global finite-energy background is Leray's
Leray--Hopf class \cite{Leray1934}.

\begin{definition}[Suitable Leray--Hopf solution on \(\R^3\times(0,T)\)]
\label{def:suitable-leray-hopf}
A suitable Leray--Hopf solution on \(\R^3\times(0,T)\) is a suitable weak
solution on \(\R^3\times(0,T)\) satisfying
\[
   u\in L^\infty(0,T;L^2(\R^3))
   \cap L^2(0,T;\dot H^1(\R^3)),
   \qquad
   u\in C_w([0,T);L^2(\R^3)),
\]
and the global Leray energy inequality.  The pressure is taken in
\(L^{3/2}_{\loc}(\R^3\times(0,T))\), modulo functions of time.
\end{definition}

Define
\begin{equation}\label{eq:C-D-def}
   A(u;z_0,r)=r^{-1}\esssup_{t_0-r^2<t<t_0}
   \int_{B_r(x_0)}|u(x,t)|^2\,dx,
\end{equation}
\begin{equation}\label{eq:C-D-def-2}
   C(u;z_0,r)=r^{-2}\iint_{Q_r(z_0)}|u|^3\,dx\,dt,
   \qquad
   D(p;z_0,r)=r^{-2}\iint_{Q_r(z_0)}
   |p-(p)_{B_r(x_0)}(t)|^{3/2}\,dx\,dt,
\end{equation}
\begin{equation}\label{eq:E-I-def}
   E(u;z_0,r)=r^{-1}\iint_{Q_r(z_0)}|\nabla u|^2\,dx\,dt,
\end{equation}
\begin{definition}[Paper--0-admissible finite-energy Type I source]\label{def:source}
A triple \((u,p,z_*)\), with \(z_*=(x_*,T)\), is a Paper--0-admissible
finite-energy Type I source if:
\begin{enumerate}[label=(\roman*)]
\item \((u,p)\) is a suitable Leray--Hopf solution on
\(\R^3\times(0,T)\);
\item \(z_*=(x_*,T)\) is a singular point;
\item the Type I bound \eqref{eq:type-I-bound} holds;
\item there are scales \(\lambda_k\downarrow0\) and a constant
\(\eta_v>0\) such that
\[
   C(u;z_*,\lambda_k)\ge\eta_v
   \qquad\hbox{for every }k;
\]
\item for the corresponding rescaled velocities
\[
   u_k(x,t)=\lambda_k u(x_*+\lambda_kx,T+\lambda_k^2t),
\]
the velocity no-escape condition holds:
\begin{equation}\label{eq:velocity-no-escape}
   \lim_{\sigma\downarrow0}
   \sup_k
   \int_{-\sigma}^{0}\int_{B_1}|u_k(x,t)|^3\,dx\,dt=0 .
\end{equation}
\end{enumerate}
The sequence \(\{\lambda_k\}\) is called an admissible Paper--0 concentration
sequence for the source.
\end{definition}

\begin{remark}
The velocity concentration in item (iv) is the velocity-only concentration
provided by the local concentration paper.  The no-escape condition in
\eqref{eq:velocity-no-escape} is the only terminal-slice compactness
assumption used to prove that the ancient limit is nonzero.  No uniform bound
on \(A+C+D+E\) over all terminal scales is assumed.
\end{remark}

\begin{lemma}[Selected kinetic energy implies no escape]
\label{lem:A-seq-no-escape}
Let \((u,p)\) satisfy the Type I bound \eqref{eq:type-I-bound}, and let
\(\lambda_k\downarrow0\).  Suppose that
\[
   \sup_k A(u;z_*,\lambda_k)\le B<\infty .
\]
Then the corresponding rescaled velocities satisfy the no-escape condition
\eqref{eq:velocity-no-escape}.
\end{lemma}

\begin{proof}
Choose \(M<\infty\) and \(t_0<T\) such that
\[
   \|u(t)\|_{L^\infty(\R^3)}
   \le M(T-t)^{-1/2},
   \qquad t_0<t<T.
\]
For \(k\) sufficiently large, \(T+\lambda_k^2t>t_0\) whenever
\(-1<t<0\).  Hence
\[
   \|u_k(t)\|_{L^\infty(B_1)}
   \le M(-t)^{-1/2},\qquad -1<t<0.
\]
By scale invariance of \(A\),
\[
   \esssup_{-1<t<0}\int_{B_1}|u_k(x,t)|^2\,dx
   =
   A(u;z_*,\lambda_k)\le B.
\]
Therefore, for \(0<\sigma<1\),
\[
\begin{aligned}
   \int_{-\sigma}^{0}\int_{B_1}|u_k|^3\,dx\,dt
   &\le
   \int_{-\sigma}^{0}
   \|u_k(t)\|_{L^\infty(B_1)}
   \int_{B_1}|u_k(x,t)|^2\,dx\,dt  \\
   &\le MB\int_{-\sigma}^{0}(-t)^{-1/2}\,dt
    =2MB\sigma^{1/2}.
\end{aligned}
\]
Taking the supremum in \(k\) and then letting \(\sigma\downarrow0\) gives
\eqref{eq:velocity-no-escape} for the tail of the sequence.  For the finitely
many remaining indices, absolute continuity of the \(L^3\)-integral on
\(B_1\times(-1,0)\) gives
\[
   \lim_{\sigma\downarrow0}
   \max_{1\le k<k_0}
   \int_{-\sigma}^{0}\int_{B_1}|u_k|^3\,dx\,dt=0.
\]
Combining the finite initial segment with the tail proves the full supremum
in \eqref{eq:velocity-no-escape}.
\end{proof}

\begin{theorem}[Epsilon regularity]\label{thm:eps-reg}
There exists \(\eps_*>0\) with the following property.  If \((u,p)\) is a
suitable weak solution in \(Q_r(z_0)\) and
\[
   C(u;z_0,r)+D(p;z_0,r)\le \eps_*,
\]
then \(u\) is bounded, and hence smooth, in \(Q_{r/2}(z_0)\).
\end{theorem}

This is the Caffarelli--Kohn--Nirenberg regularity criterion
\cite{CKN1982}; see also Lin's direct proof \cite{Lin1998}.  We use its
contrapositive in the concentration argument below.

\begin{proposition}[Paper--0 velocity concentration input]
\label{prop:paper0-velocity-concentration}
Let \((u,p)\) be suitable in a parabolic neighborhood of
\(z_*=(x_*,T)\).  If \(z_*\) is singular, then there exist
\(\lambda_k\downarrow0\) and \(\eta_v>0\) such that
\[
   C(u;z_*,\lambda_k)\ge\eta_v
   \qquad\hbox{for every }k.
\]
\end{proposition}

\begin{proof}
This is the velocity-only concentration corollary of the local concentration
paper \cite{Paper0LocalConcentration}.  It is obtained there from the
velocity epsilon-regularity criterion: if
\(\limsup_{r\downarrow0}C(u;z_*,r)=0\), then \(u\) is bounded in a backward
parabolic cylinder ending at \(z_*\), contradicting singularity of \(z_*\).
\end{proof}

\begin{lemma}[Scaling]\label{lem:scaling}
Let \((u,p)\) be suitable in a neighborhood of \((x_*,T)\).  For
\(\lambda>0\), set
\[
   u^\lambda(x,t)=\lambda u(x_*+\lambda x,T+\lambda^2t),
   \qquad
   p^\lambda(x,t)=\lambda^2p(x_*+\lambda x,T+\lambda^2t).
\]
Then \((u^\lambda,p^\lambda)\) is suitable on the rescaled domain.  Moreover
\(A,C,D,E\) are invariant under this scaling.
\end{lemma}

\begin{proof}
Suitability of the rescaled pair and preservation of pressure gauges follow
by applying \eqref{eq:lei} to the test function
\[
   \phi^\lambda(X,S)
   =
   \phi\left(\frac{X-x_*}{\lambda},\frac{S-T}{\lambda^2}\right)
\]
and then changing variables
\((X,S)=(x_*+\lambda x,T+\lambda^2t)\).  The differential identities
\[
   \partial_tu^\lambda=\lambda^3(\partial_tu)(X,S),\quad
   (u^\lambda\cdot\nabla)u^\lambda
      =\lambda^3((u\cdot\nabla)u)(X,S),
\]
\[
   \nabla p^\lambda=\lambda^3(\nabla p)(X,S),\qquad
   \Delta u^\lambda=\lambda^3(\Delta u)(X,S)
\]
give the rescaled equation and the divergence condition.

For the pressure average,
\[
   (p^\lambda)_{B_r}(t)
   =
   \lambda^2(p)_{B_{\lambda r}(x_*)}(S).
\]
Consequently
\[
   \iint_{Q_r}|u^\lambda|^3\,dx\,dt
   =
   \lambda^{-2}\iint_{Q_{\lambda r}(z_*)}|u|^3\,dX\,dS
\]
and
\[
   \iint_{Q_r}|p^\lambda-(p^\lambda)_{B_r}(t)|^{3/2}\,dx\,dt
   =
   \lambda^{-2}
   \iint_{Q_{\lambda r}(z_*)}
   |p-(p)_{B_{\lambda r}(x_*)}(S)|^{3/2}\,dX\,dS .
\]
Multiplication by \(r^{-2}\) therefore gives invariance of \(C\) and \(D\).
The remaining two quantities have the same critical dimensional balance:
\begin{align*}
   \int_{B_r}|u^\lambda(x,t)|^2\,dx
   &=
   \lambda^{-1}\int_{B_{\lambda r}(x_*)}|u(X,S)|^2\,dX,\\
   \iint_{Q_r}|\nabla u^\lambda|^2\,dx\,dt
   &=
   \lambda^{-1}\iint_{Q_{\lambda r}(z_*)}|\nabla u|^2\,dX\,dS .
\end{align*}
Multiplication by \(r^{-1}\) on the rescaled cylinder is the same as
multiplication by \((\lambda r)^{-1}\) on the original cylinder.  Hence
\(A\) and \(E\) are invariant as well.
\end{proof}

\begin{lemma}[Pressure gauges do not change the equations]
\label{lem:gauges}
Let \((v,q)\) be suitable in a cylinder \(Q\), and let
\(\alpha(t)\in L^{3/2}_{\loc}\) depend only on time.  Then
\((v,q-\alpha(t))\) satisfies the same distributional Navier--Stokes equation,
the same divergence condition, and the same local energy inequality.
\end{lemma}

\begin{proof}
In the distributional equation, replacing \(q\) by \(q-\alpha(t)\) adds
\(-\nabla\alpha(t)=0\).  In the local energy inequality the only new term is
\[
   -2\int \alpha(t)v(x,t)\cdot\nabla\phi(x,t)\,dx\,dt .
\]
The product is integrable on the support of \(\phi\): by suitability,
\(v\in L^\infty_tL^2_x\) locally, and
\[
   \left|\int v(\cdot,t)\cdot\nabla\phi(\cdot,t)\,dx\right|
   \le
   \|v(t)\|_{L^2(\supp\phi(t))}
   \|\nabla\phi(t)\|_{L^2}.
\]
For almost every \(t\), \(\divergence v(\cdot,t)=0\) in distributions and
\(\phi(\cdot,t)\) is compactly supported, so
\[
   \int v(x,t)\cdot\nabla\phi(x,t)\,dx=0 .
\]
Thus the additional term vanishes.
\end{proof}

\begin{lemma}[Pressure-normalized concentration at singular points]
\label{lem:singular-concentration}
If \(z_*=(x_*,T)\) is singular, then there are \(\eta>0\) and \(r_0>0\) such
that, for every \(0<r<r_0\),
\begin{equation}\label{eq:positive-concentration}
   C(u;z_*,r)+D(p;z_*,r)\ge \eta .
\end{equation}
\end{lemma}

\begin{proof}
Choose \(r_0>0\) so that \(Q_r(z_*)\) lies in the domain for every
\(0<r<r_0\), and take \(\eta=\eps_*\).  If for some \(0<r<r_0\)
\[
   C(u;z_*,r)+D(p;z_*,r)\le\eps_*,
\]
then \cref{thm:eps-reg} would imply that \(u\) is bounded in
\(Q_{r/2}(z_*)\).  This is exactly regularity of \(z_*\) in
\cref{def:terminal-singular}, contradicting singularity.  Hence the strict
lower bound holds at every such radius.
\end{proof}

\section{The Type I blow-up sequence and compactness}
\label{sec:blowup-compactness}

\subsection{The rescaled sequence}

Choose \(M<\infty\) so that, after decreasing \(T_0<T\) if necessary,
\begin{equation}\label{eq:type-I-pointwise}
   \|u(t)\|_{L^\infty(\R^3)}
   \le \frac{M}{\sqrt{T-t}},
   \qquad T_0<t<T.
\end{equation}
For \(\lambda_k\downarrow0\), define
\begin{equation}\label{eq:rescaled-sequence}
   u_k(x,t)=\lambda_k u(x_*+\lambda_k x,T+\lambda_k^2t),
   \qquad
   p_k(x,t)=\lambda_k^2p(x_*+\lambda_k x,T+\lambda_k^2t).
\end{equation}
Then \(u_k\) is defined on \(\R^3\times(-T/\lambda_k^2,0)\), and these
domains exhaust \(\R^3\times(-\infty,0)\).

\subsection{Inherited Type I bounds}

\begin{lemma}[Inherited Type I bound]\label{lem:inherited-bound}
For every compact interval \(I\Subset(-\infty,0)\),
\[
   \sup_k\|u_k\|_{L^\infty(\R^3\times I)}<\infty.
\]
More precisely, for every fixed \(t<0\) and all sufficiently large \(k\),
\[
   \|u_k(t)\|_{L^\infty(\R^3)}
   \le \frac{M}{\sqrt{-t}}.
\]
\end{lemma}

\begin{proof}
Since \(T-(T+\lambda_k^2t)=\lambda_k^2(-t)\), \eqref{eq:type-I-pointwise}
gives, for all large \(k\),
\[
   \|u_k(t)\|_\infty
   =
   \lambda_k\|u(T+\lambda_k^2t)\|_\infty
   \le
   \lambda_k\frac{M}{\sqrt{\lambda_k^2(-t)}}
   =
   \frac{M}{\sqrt{-t}}.
\]
Taking the supremum over \(t\in I\) gives the interval statement.
\end{proof}

\subsection{Local energy, pressure, and time-derivative estimates}

\begin{lemma}[Pressure gauges and local pressure bounds]\label{lem:pressure}
After fixing the whole-space Leray pressure representative below, set
\[
   a_k(t)=(p_k)_{B_1}(t)
\]
whenever the rescaled solution is defined.  Then for every
\[
   Q=B_R\times[-S,-\sigma]\Subset\R^3\times(-\infty,0)
\]
one has
\[
   \sup_k\|p_k-a_k(t)\|_{L^{3/2}(Q)}<\infty .
\]
\end{lemma}

\begin{proof}
For large \(k\), the original times \(T+\lambda_k^2t\), \(t\in[-S,-\sigma]\),
lie in the interval where \eqref{eq:type-I-pointwise} holds.  On those slices
\(u(t)\in L^2(\R^3)\cap L^\infty(\R^3)\), hence
\[
   u_i(t)u_j(t)\in L^1(\R^3)\cap L^{3/2}(\R^3).
\]
Use the whole-space Leray pressure representative
\[
   p^L=\mathcal R_i\mathcal R_j(u_i u_j),
\]
modulo functions of time.  We use the standard convention
\[
   \widehat{\mathcal R_i f}(\xi)=i\xi_i|\xi|^{-1}\widehat f(\xi),
\]
so that \(\mathcal R_i\mathcal R_j\) has multiplier
\(-\xi_i\xi_j|\xi|^{-2}\) and
\(-\Delta(\mathcal R_i\mathcal R_j f)=\partial_i\partial_jf\).  Thus this is
the pressure reconstruction from
\(-\Delta p=\partial_i\partial_j(u_i u_j)\) and the Calderon--Zygmund theory
of Riesz transforms \cite{Stein1970,Sohr2001}.  Since
\(u(t)\in L^2(\R^3)\cap L^\infty(\R^3)\), interpolation gives
\(u(t)\in L^3(\R^3)\), hence \(u_i(t)u_j(t)\in L^{3/2}(\R^3)\) and
\(p^L(t)\in L^{3/2}(\R^3)\) for a.e. \(t\).

The given pressure and \(p^L\) differ by a function of time.  Indeed, applying
the Helmholtz projection to the whole-space weak momentum equation gives
\[
   \partial_tu-\Delta u+\mathbb P\divergence(u\otimes u)=0
\]
in solenoidal distributions.  The complementary gradient part of
\(\divergence(u\otimes u)\) is precisely
\(-\nabla p^L\) with the convention above.  Therefore the same velocity field
solves the distributional equation with \(p^L\) in place of \(p\).  Subtracting
the two formulations gives
\[
   \nabla(p-p^L)=0
   \quad\hbox{in }\mathcal D'(\R^3)
\]
for a.e. time.  By the de Rham theorem on \(\R^3\), \(p-p^L\) is spatially
constant for a.e. time.  Replacing the pressure by \(p^L\) is therefore a
time-dependent gauge change and does not affect suitability by
\cref{lem:gauges}.

Fix \(\chi_R\in C_c^\infty(B_{8R})\) with \(\chi_R=1\) on \(B_{4R}\).  For
\(t\in[-S,-\sigma]\), choose the representative so that, in \(B_{2R}\),
\[
   p_k=q_k+h_k,\qquad
   q_k=\mathcal R_i\mathcal R_j(\chi_Ru_{k,i}u_{k,j}),
\]
and \(h_k(\cdot,t)\) is harmonic in \(B_{4R}\) for a.e. \(t\).  This choice
is made before \(a_k(t)=(p_k)_{B_1}(t)\) is evaluated.  The harmonicity holds
because \((1-\chi_R)u_{k,i}u_{k,j}\) is supported outside \(B_{4R}\), so the
localized Riesz potential generated by this exterior part has zero Laplacian
in \(B_{4R}\).  Calderon--Zygmund boundedness gives, for a.e. \(t\),
\[
   \|q_k(t)\|_{L^{3/2}(B_{2R})}
   \le C\|u_k(t)\|_{L^3(B_{8R})}^2
   \le C(R,\sigma,M).
\]
Integrating over \([-S,-\sigma]\) gives a uniform
\(L^{3/2}_{x,t}\)-bound for \(q_k\).
Let \(c_{k,R}(t)=h_k(0,t)\), where \(h_k(\cdot,t)\) is represented by its
smooth harmonic representative.  For \(x\in B_R\), the kernel estimate
\[
   |K_{ij}(x-z)-K_{ij}(-z)|\le C R|z|^{-4},\qquad |z|>4R,
\]
implies
\[
   |h_k(x,t)-c_{k,R}(t)|
   \le C R\sum_{m=2}^\infty(2^mR)^{-4}
   \int_{2^mR<|z|<2^{m+1}R}|u_k(z,t)|^2\,dz .
\]
Using \cref{lem:inherited-bound},
\[
   \int_{2^mR<|z|<2^{m+1}R}|u_k(z,t)|^2\,dz
   \le C M^2\sigma^{-1}(2^mR)^3,
\]
and therefore
\[
   |h_k(x,t)-c_{k,R}(t)|
   \le C M^2\sigma^{-1}\sum_{m=2}^{\infty}2^{-m},
   \qquad x\in B_R,
\]
which implies
\[
   \|h_k(t)-c_{k,R}(t)\|_{L^\infty(B_R)}\le C(M,\sigma).
\]
Thus
\[
   \sup_k\|p_k-c_{k,R}(t)\|_{L^{3/2}(B_R\times[-S,-\sigma])}<\infty .
\]
If \(R<1\), enlarge \(R\) to \(1\).  Otherwise \(B_1\subset B_R\).  Jensen's
inequality gives, for a.e. \(t\),
\[
   |a_k(t)-c_{k,R}(t)|^{3/2}
   \le
   |B_1|^{-1}\int_{B_1}|p_k-c_{k,R}(t)|^{3/2}\,dx .
\]
Integrating in time and adding constants on \(B_R\) gives the desired bound
for \(p_k-a_k(t)\).
\end{proof}

\begin{lemma}[Compatibility of pressure gauges]\label{lem:gauge-compatibility}
Let \(B_1\subset B_R\), \(I\Subset(-\infty,0)\), and suppose
\[
   p_k-c_k(t)
   \quad\hbox{is bounded in }L^{3/2}(B_R\times I)
\]
for some time-dependent gauges \(c_k(t)\).  If
\(a_k(t)=(p_k)_{B_1}(t)\), then \(p_k-a_k(t)\) is bounded in
\(L^{3/2}(B_{R'}\times I)\) for every \(R'\le R\).  Changing between such
gauges changes the pressure only by a function of time and therefore does not
affect the equation, the local energy inequality, or the centered equation.
\end{lemma}

\begin{proof}
For a.e. \(t\), Jensen's inequality gives
\[
   |a_k(t)-c_k(t)|^{3/2}
   \le |B_1|^{-1}\int_{B_1}|p_k-c_k(t)|^{3/2}\,dx .
\]
Integrating in \(t\) and using \(B_{R'}\subset B_R\) yields
\[
   \|a_k-c_k\|_{L^{3/2}(I)}^{3/2}
   \le
   |B_1|^{-1}\|p_k-c_k(t)\|_{L^{3/2}(B_R\times I)}^{3/2}.
\]
Thus
\[
   \|p_k-a_k(t)\|_{L^{3/2}(B_{R'}\times I)}
   \le
   \|p_k-c_k(t)\|_{L^{3/2}(B_{R'}\times I)}
   +|B_{R'}|^{2/3}\|a_k-c_k\|_{L^{3/2}(I)}
\]
and the right-hand side is bounded by
\[
   C(R')\|p_k-c_k(t)\|_{L^{3/2}(B_R\times I)} .
\]
This proves the asserted bound.
The last assertion is \cref{lem:gauges}, and centering transforms the
time-dependent gauge into another time-dependent gauge.
\end{proof}

\begin{lemma}[Local energy and dissipation]\label{lem:energy-diss}
Let
\[
   Q=B_R\times[-S,-\sigma]\Subset\R^3\times(-\infty,0).
\]
Then
\[
   \sup_k\|u_k\|_{L^\infty_tL^2_x(Q)}
   +\sup_k\|\nabla u_k\|_{L^2(Q)}<\infty .
\]
\end{lemma}

\begin{proof}
The \(L^\infty_tL^2_x\) estimate follows from
\[
   \|u_k(t)\|_{L^2(B_R)}
   \le |B_R|^{1/2}M\sigma^{-1/2},
   \qquad t\in[-S,-\sigma].
\]
Choose \(\phi\in C_c^\infty(B_{2R}\times[-2S,-\sigma/2])\), \(0\le\phi\le1\),
with \(\phi=1\) on \(Q\).  Apply the local energy inequality to
\((u_k,p_k-a_k(t))\), using \cref{lem:pressure} on the larger cylinder.  For
a.e. admissible times and then by monotone approximation of temporal cutoffs,
\[
   2\iint |\nabla u_k|^2\phi
   \le
   \iint |u_k|^2|\partial_t\phi+\Delta\phi|
   +\iint (|u_k|^2+2|p_k-a_k(t)|)|u_k||\nabla\phi|.
\]
The first term is bounded because \(\partial_t\phi+\Delta\phi\) is bounded
and \(u_k\) is uniformly bounded in \(L^\infty\) on the finite-measure support
of \(\phi\).  The cubic velocity term satisfies
\[
   \iint |u_k|^3|\nabla\phi|
   \le
   \|\nabla\phi\|_\infty
   |\supp\nabla\phi|\,\|u_k\|_{L^\infty(\supp\nabla\phi)}^3
   \le C(Q,\phi,M,\sigma).
\]
For the pressure term, H\"older's inequality on \(\supp\nabla\phi\) gives
\[
   \iint |p_k-a_k(t)|\,|u_k|\,|\nabla\phi|
   \le
   \|p_k-a_k(t)\|_{L^{3/2}}
   \|u_k\nabla\phi\|_{L^3},
\]
and both factors are uniformly bounded by \cref{lem:pressure} and
\cref{lem:inherited-bound}.  Since \(\phi=1\) on \(Q\), the dissipation bound
follows.
\end{proof}

\begin{lemma}[Time derivative bound]\label{lem:time-derivative}
For every compact cylinder \(Q\Subset\R^3\times(-\infty,0)\),
\[
   \partial_tu_k
   \quad\hbox{is bounded in }L^{3/2}_tW^{-1,3/2}_x(Q).
\]
\end{lemma}

\begin{proof}
The equation is
\[
   \partial_tu_k
   =
   \Delta u_k-\divergence(u_k\otimes u_k)-\nabla(p_k-a_k(t)).
\]
Let \(Q=B_R\times[-S,-\sigma]\).  The term \(\Delta u_k\) is bounded in
\(L^2_tH^{-1}_x(Q)\) by \cref{lem:energy-diss}; since
\(W^{1,3}_0(B_R)\subset H^1_0(B_R)\) and the time interval is finite, this is
also an \(L^{3/2}_tW^{-1,3/2}_x\)-bound.  Explicitly, for
\(\zeta\in W^{1,3}_0(B_R)\),
\[
   |\langle\Delta u_k,\zeta\rangle|
   =
   \left|\int_{B_R}\nabla u_k:\nabla\zeta\,dx\right|
   \le
   C(R)\|\nabla u_k\|_{L^2(B_R)}\|\zeta\|_{W^{1,3}(B_R)}.
\]
The \(L^2_t\)-bound on \(\nabla u_k\), together with the finite length of
\([-S,-\sigma]\), gives the claimed \(L^{3/2}_t\)-bound.

For the nonlinear term, for
\(\zeta\in W^{1,3}_0(B_R)\),
\[
   \left|\int_{B_R}(u_k\otimes u_k):\nabla\zeta\,dx\right|
   \le
   \|u_k\otimes u_k\|_{L^{3/2}(B_R)}
   \|\nabla\zeta\|_{L^3(B_R)}.
\]
Moreover
\[
   \|u_k\otimes u_k\|_{L^{3/2}(Q)}
   \le
   \|u_k\|_{L^\infty(Q)}\|u_k\|_{L^{3/2}(Q)}
   \le C(Q).
\]
Thus \(\divergence(u_k\otimes u_k)\) is bounded in
\(L^{3/2}_tW^{-1,3/2}_x(Q)\).
For the pressure term,
\[
   \left|\int_{B_R}(p_k-a_k(t))\,\divergence\zeta\,dx\right|
   \le
   \|p_k-a_k(t)\|_{L^{3/2}(B_R)}
   \|\nabla\zeta\|_{L^3(B_R)}.
\]
\Cref{lem:pressure} gives
\[
   p_k-a_k(t)\quad\hbox{bounded in }L^{3/2}(Q).
\]
Hence \(\nabla(p_k-a_k(t))\) is bounded in
\(L^{3/2}_tW^{-1,3/2}_x(Q)\).  Summing the three terms gives the result.
\end{proof}

\subsection{Compactness away from the terminal time}

\begin{proposition}[Compactness]\label{prop:compactness}
After passing to a subsequence,
\[
   u_k\to U
\]
strongly in \(L^q_{\loc}(\R^3\times(-\infty,0))\) for every finite \(q\), and
\[
   p_k-a_k(t)\rightharpoonup Q
   \quad\hbox{weakly in }L^{3/2}_{\loc}(\R^3\times(-\infty,0)).
\]
\end{proposition}

\begin{proof}
Fix \(Q_0\Subset\R^3\times(-\infty,0)\).  By
\cref{lem:energy-diss}, \(u_k\) is bounded in \(L^2_tH^1_x(Q_0)\), and by
\cref{lem:time-derivative}, \(\partial_tu_k\) is bounded in
\(L^{3/2}_tW^{-1,3/2}_x(Q_0)\).  On bounded balls,
\[
   H^1\Subset L^2\hookrightarrow W^{-1,3/2}.
\]
The Aubin--Lions--Simon theorem \cite{Aubin1963,LionsMagenes,Simon1987} gives
relative compactness in \(L^2(Q_0)\).  Hence, after passing to a subsequence,
\(u_k\to U\) strongly in \(L^2(Q_0)\) and a.e. on \(Q_0\).

The inherited Type I bound gives a common \(L^\infty(Q_0)\) bound, and the
limit belongs to the same \(L^\infty\) ball by a.e. convergence.  For
\(q\ge2\),
\[
   \|u_k-U\|_{L^q(Q_0)}^q
   \le
   (2\sup_j\|u_j\|_{L^\infty(Q_0)}+2\|U\|_{L^\infty(Q_0)})^{q-2}
   \|u_k-U\|_{L^2(Q_0)}^2.
\]
For \(1\le q<2\), H\"older's inequality on the finite-measure set \(Q_0\)
gives
\[
   \|u_k-U\|_{L^q(Q_0)}
   \le |Q_0|^{1/q-1/2}\|u_k-U\|_{L^2(Q_0)}\to0.
\]
Choose an exhaustion \(Q^{(m)}\Subset\R^3\times(-\infty,0)\) by compact
cylinders.  The preceding argument gives a subsequence converging on
\(Q^{(1)}\); extracting successively on \(Q^{(m)}\) and taking the diagonal
subsequence gives strong local convergence on every compact cylinder.

For pressure, \cref{lem:pressure} gives boundedness of
\(p_k-a_k(t)\) in \(L^{3/2}(Q^{(m)})\) for every \(m\).  Since
\(L^{3/2}\) is reflexive, each restriction has a weakly convergent
subsequence.  The same diagonal extraction gives a single limit \(Q\) in
\(L^{3/2}_{\loc}\), because the weak limits agree on overlaps by uniqueness
of distributional limits.
\end{proof}

\subsection{The ancient suitable weak limit}

\begin{lemma}[Limit equation]\label{lem:limit-equation}
The limit \((U,Q)\) solves
\[
   \partial_tU+(U\cdot\nabla)U+\nabla Q=\Delta U,\qquad
   \divergence U=0
\]
in distributions on \(\R^3\times(-\infty,0)\).
\end{lemma}

\begin{proof}
For \(\varphi\in C_c^\infty(\R^3\times(-\infty,0);\R^3)\), the weak
formulation for \((u_k,p_k-a_k(t))\) is
\[
   \iint\left[
      -u_k\cdot\partial_t\varphi
      -(u_k\otimes u_k):\nabla\varphi
      +(p_k-a_k(t))\divergence\varphi
      +\nabla u_k:\nabla\varphi
   \right]\,dx\,dt=0.
\]
The linear time term passes by strong \(L^2_{\loc}\)-convergence of \(u_k\).
The pressure term passes by weak \(L^{3/2}_{\loc}\)-convergence of
\(p_k-a_k(t)\), since \(\divergence\varphi\in L^3\).  The viscous term passes
by weak \(L^2_{\loc}\)-convergence of \(\nabla u_k\), which follows from the
uniform local dissipation bound.  The nonlinear term
passes because
\[
   \|u_k\otimes u_k-U\otimes U\|_{L^1(\supp\varphi)}
   \le
   \|u_k-U\|_{L^2}
   \bigl(\|u_k\|_{L^2}+\|U\|_{L^2}\bigr)\to0.
\]
Testing with gradients of scalar functions gives \(\divergence U=0\).
\end{proof}

\begin{lemma}[Stability of suitability]\label{lem:suitability}
The pair \((U,Q)\) is a suitable weak solution on
\(\R^3\times(-\infty,0)\).
\end{lemma}

\begin{proof}
Use the distributional form of the local energy inequality.  For every
non-negative \(\phi\in C_c^\infty(\R^3\times(-\infty,0))\),
\[
   2\iint |\nabla u_k|^2\phi
   \le
   \iint |u_k|^2(\partial_t\phi+\Delta\phi)
   +\iint (|u_k|^2+2P_k)u_k\cdot\nabla\phi ,
   \qquad P_k=p_k-a_k(t).
\]
This follows from \eqref{eq:lei} by applying \eqref{eq:lei} with temporal
cutoffs which increase to \(1\) on the time projection of \(\supp\phi\).  The
monotone convergence theorem passes the nonnegative dissipation term to the
displayed distributional form.

The terms with \(|u_k|^2\) converge strongly in \(L^1\) because
\[
   \||u_k|^2-|U|^2\|_{L^1(\supp\phi)}
   \le
   \|u_k-U\|_{L^2}
   \bigl(\|u_k\|_{L^2}+\|U\|_{L^2}\bigr)\to0.
\]
The cubic flux converges strongly since \(u_k\to U\) strongly in
\(L^3_{\loc}\), hence
\[
   |u_k|^2u_k\to |U|^2U
   \quad\hbox{in }L^1_{\loc}.
\]
For the pressure flux, \(u_k\cdot\nabla\phi\to U\cdot\nabla\phi\) strongly in
\(L^3\), while \(P_k\rightharpoonup Q\) weakly in \(L^{3/2}\).  Therefore
\[
   \iint P_k u_k\cdot\nabla\phi\to
   \iint Q\,U\cdot\nabla\phi .
\]
Finally,
\(\nabla u_k\rightharpoonup\nabla U\) in \(L^2_{\loc}\), and
\[
   w\mapsto\iint |w|^2\phi
\]
is weakly lower semicontinuous because \(\phi\ge0\) and the expression equals
\(\|\sqrt{\phi}\,w\|_{L^2}^2\).  Passing to the limit gives the local energy
inequality for \((U,Q)\) in distributional form.  The endpoint form in
\cref{def:suitable} is recovered by applying the distributional inequality to
standard time cutoffs approximating characteristic functions of
\((t_1,t_2)\); weak \(L^2_{\loc}\)-continuity of local energy-class solutions
gives the endpoint traces for a.e. \(t_1<t_2\).
\end{proof}

\begin{lemma}[Inherited ancient Type I bound]\label{lem:limit-bound}
For every compact interval \(I=[-S,-\sigma]\Subset(-\infty,0)\),
\[
   \|U\|_{L^\infty(\R^3\times I)}\le M\sigma^{-1/2}.
\]
\end{lemma}

\begin{proof}
\Cref{lem:inherited-bound} gives
\[
   \|u_k\|_{L^\infty(\R^3\times I)}\le M\sigma^{-1/2}.
\]
Fix \(R<\infty\).  Since \(u_k\to U\) a.e. on \(B_R\times I\), the essential
supremum of \(U\) on \(B_R\times I\) cannot exceed \(M\sigma^{-1/2}\).  If it
did, \(|u_k|\) would exceed this bound by a positive amount on a positive
measure subset for large \(k\).  Letting \(R\to\infty\) through integers gives
the result on \(\R^3\times I\).
\end{proof}

\begin{proposition}[Ancient suitable Type I limit]\label{prop:ancient-limit}
The pair \((U,Q)\) is an ancient suitable weak solution on
\(\R^3\times(-\infty,0)\) satisfying the interval Type I bound in
\cref{lem:limit-bound}.
\end{proposition}

\begin{proof}
Combine \cref{lem:limit-equation,lem:suitability,lem:limit-bound}.
\end{proof}

\section{Persistence of concentration and nontriviality}
\label{sec:nontriviality}

\subsection{Persistence of terminal concentration away from the terminal slice}

The singularity of the source enters the proof through the velocity
concentration sequence in \cref{def:source}, which is supplied by
\cref{prop:paper0-velocity-concentration}.  The no-escape condition is then
exactly the assertion that this velocity mass does not concentrate entirely in
the terminal layer \(t=0\) after rescaling.

\subsection{Scale selection}

\begin{lemma}[Scale selection from local concentration]\label{lem:scale-selection}
Let \(\lambda_k\downarrow0\) be an admissible Paper--0 concentration sequence
for the source in the sense of \cref{def:source}.  Then there are
\(\sigma\in(0,1)\), a compact cylinder
\[
   K_0=B_1\times(-1,-\sigma)\Subset\R^3\times(-\infty,0),
\]
and \(\eta_1>0\) such that
\begin{equation}\label{eq:selected-lower-bound}
   \iint_{K_0}|u_k|^3\,dx\,dt\ge\eta_1
\end{equation}
for all \(k\).
\end{lemma}

\begin{proof}
By scale invariance of \(C\),
\[
   \int_{-1}^{0}\int_{B_1}|u_k|^3\,dx\,dt
   =
   C(u;z_*,\lambda_k)\ge\eta_v .
\]
The no-escape condition \eqref{eq:velocity-no-escape} gives
\(\sigma\in(0,1)\) such that
\[
   \sup_k\int_{-\sigma}^{0}\int_{B_1}|u_k|^3\,dx\,dt
   \le \frac{\eta_v}{2}.
\]
Therefore
\[
   \int_{-1}^{-\sigma}\int_{B_1}|u_k|^3\,dx\,dt
   \ge \frac{\eta_v}{2}
\]
for every \(k\).  The assertion follows with
\(\eta_1=\eta_v/2\).
\end{proof}

\subsection{Nonzero ancient limits}

\begin{proposition}[Nonzero ancient limit]\label{prop:nonzero}
For the scales chosen in \cref{lem:scale-selection}, the limit \(U\) in
\cref{prop:ancient-limit} is nonzero.  More precisely, there are a compact
cylinder \(K_0\Subset\R^3\times(-\infty,0)\) and \(\eta_1>0\) such that
\[
   \iint_{K_0}|U|^3\,dx\,dt\ge\eta_1 .
\]
\end{proposition}

\begin{proof}
Let \(K_0=B_1\times(-1,-\sigma)\) and \(\eta_1>0\) be given by
\cref{lem:scale-selection}.  Since
\(K_0\Subset\R^3\times(-\infty,0)\), the compactness theorem gives, after
passing to the subsequence used to define \(U\),
\[
   u_k\to U\qquad\hbox{strongly in }L^3(K_0).
\]
Hence
\[
   \int_{K_0}|U|^3\,dx\,dt
   =
   \lim_{k\to\infty}\int_{K_0}|u_k|^3\,dx\,dt
   \ge\eta_1 .
\]
\end{proof}

The preceding proposition is the point at which singularity of the source
enters the compactness argument.  Without the velocity concentration and the
no-escape condition, the compactness theorem alone could produce the zero
ancient solution.

\section{Centered variables and normalized Seregin limits}\label{sec:seregin}

\subsection{The centered self-similar equation}

Set
\[
   y=\frac{x}{\sqrt{-t}},\qquad
   \tau=-\log(-t),\qquad
   V(y,\tau)=\sqrt{-t}\,U(y\sqrt{-t},t),
\]
and
\[
   \Pi(y,\tau)=(-t)Q(y\sqrt{-t},t),
\]
up to addition of a function of \(\tau\).

\begin{lemma}[Renormalized equation]\label{lem:renormalized-equation}
The pair \((V,\Pi)\) solves \eqref{eq:centered}.
\end{lemma}

\begin{proof}
Write \(s=-t=e^{-\tau}\), \(x=s^{1/2}y\), and
\[
   U(x,t)=s^{-1/2}V(y,\tau).
\]
The identities are first obtained in distributions by testing the equation
for \(U\) against the inverse change of variables.  After
\cref{lem:smoothness}, the same identities hold classically for the smooth
representative.  Because \(s_t=-1\),
\(\tau_t=s^{-1}\), and \(y_t=(2s)^{-1}y\),
\[
   \partial_tU
   =
   s^{-3/2}\left(\partial_\tau V+\frac12V+\frac12y\cdot\nabla V\right),
\]
\[
   (U\cdot\nabla_x)U=s^{-3/2}(V\cdot\nabla)V,\qquad
   \Delta_xU=s^{-3/2}\Delta V,\qquad
   \nabla_xQ=s^{-3/2}\nabla\Pi.
\]
Substitution into the equation for \(U\) gives \eqref{eq:centered}, and the
divergence condition is preserved by the spatial scaling.
\end{proof}

\subsection{Boundedness, smoothness, and local energy}

\begin{lemma}[Smoothness]\label{lem:smoothness}
The ancient solution \(U\) is smooth on \(\R^3\times(-\infty,0)\), and
\((V,\Pi)\) is smooth on \(\R^3\times\R\) after fixing a pressure gauge.
\end{lemma}

\begin{proof}
On each compact cylinder \(Q_0\Subset\R^3\times(-\infty,0)\),
\cref{lem:limit-bound} gives \(U\in L^\infty(Q_0)\).  Hence
\(U\in L^s_tL^r_x(Q_0)\) for every finite \(r,s\).  Choose \(r=s=10\), so
\[
   \frac2s+\frac3r=\frac12<1.
\]
Serrin's interior regularity criterion \cite{Serrin1962,Sohr2001} applies and
gives local H\"older regularity of \(U\) on a slightly smaller cylinder.
Since \(Q_0\) was arbitrary, this removes every possible interior singular
point in \(\R^3\times(-\infty,0)\).  The bootstrap starts from
\[
   \partial_tU-\Delta U+\nabla Q=-\divergence(U\otimes U),\qquad
   \divergence U=0,
\]
with \(U\otimes U\) locally bounded after Serrin regularity.  Interior Stokes
estimates on nested cylinders give \(W^{2,1}_q\)-regularity for every finite
\(q\).  Parabolic Sobolev embedding then improves the H\"older regularity of
\(\nabla U\), and differentiating the equation and iterating gives
\[
   U\in C^\infty(\R^3\times(-\infty,0)).
\]
The pressure is recovered locally from
\(-\Delta Q=\partial_i\partial_j(U_iU_j)\) plus a harmonic function and is
smooth after fixing a time-dependent gauge.  The centered change of variables
is a smooth diffeomorphism for \(t<0\), so \(V\) and \(\Pi\) are smooth.
\end{proof}

\begin{corollary}[Classical Type I bound]\label{cor:classical-type-I-bound}
For the smooth representative from \cref{lem:smoothness},
\[
   \|U(t)\|_{L^\infty(\R^3)}\le \frac{M}{\sqrt{-t}},
   \qquad t<0.
\]
\end{corollary}

\begin{proof}
Fix \(t_0<0\).  Choose \(\delta_n\downarrow0\) with
\[
   I_n=[t_0-\delta_n,t_0+\delta_n]\Subset(-\infty,0).
\]
\Cref{lem:limit-bound} gives
\[
   \|U\|_{L^\infty(\R^3\times I_n)}
   \le M(-t_0-\delta_n)^{-1/2}.
\]
Since the representative is continuous, the essential spacetime bound holds
pointwise.  Letting \(n\to\infty\) gives the estimate at \(t_0\).
\end{proof}

\begin{lemma}[Renormalized boundedness and nontriviality]
\label{lem:renorm-bounds}
The centered field satisfies
\[
   \|V\|_{L^\infty(\R^3\times\R)}\le M.
\]
Moreover \eqref{eq:nonvanishing} holds for some compact cylinder
\(K\Subset\R^3\times\R\) and \(\eta_0>0\).
\end{lemma}

\begin{proof}
The \(L^\infty\) bound follows from \cref{cor:classical-type-I-bound}:
\[
   |V(y,\tau)|=\sqrt{-t}|U(y\sqrt{-t},t)|\le M.
\]
For nontriviality, \cref{prop:nonzero} gives a positive velocity lower bound
on a compact cylinder in \((x,t)\).  Under
\((x,t)=(e^{-\tau/2}y,-e^{-\tau})\),
\[
   dx\,dt=e^{-5\tau/2}\,dy\,d\tau,\qquad
   U=e^{\tau/2}V,
\]
and hence
\[
   |U|^3\,dx\,dt=e^{-\tau}|V|^3\,dy\,d\tau .
\]
On the compact image cylinder the factor \(e^{-\tau}\) is bounded above and
below by positive constants, so the positive lower bound becomes
\eqref{eq:nonvanishing}, with a possibly different \(\eta_0\).
\end{proof}

\begin{lemma}[Local energy bounds in centered variables]
\label{lem:centered-local-energy}
Let \(V\) be obtained from a Paper--0-admissible source.  For every
\(R<\infty\) and
compact interval \(I\subset\R\),
\[
   \int_I\int_{B_R}|V|^2\,dy\,d\tau
   +
   \int_I\int_{B_R}|\nabla V|^2\,dy\,d\tau<\infty .
\]
After subtracting a function of \(\tau\), the associated pressure satisfies
\[
   \Pi-b(\tau)\in L^{3/2}(B_R\times I).
\]
\end{lemma}

\begin{proof}
The unrenormalized limit \((U,Q)\) is suitable and has local energy,
dissipation, and pressure bounds on compact negative-time cylinders by
\cref{lem:energy-diss,lem:pressure,prop:ancient-limit}; the bounds for the
limit follow from the weak lower semicontinuity and weak compactness already
used in \cref{lem:suitability}.  On each compact
\((y,\tau)\)-cylinder, the map
\[
   (y,\tau)\mapsto (x,t)=(e^{-\tau/2}y,-e^{-\tau})
\]
has Jacobian \(dx\,dt=e^{-5\tau/2}\,dy\,d\tau\), with
\[
   V=e^{-\tau/2}U,\qquad
   \nabla_yV=e^{-\tau}\nabla_xU,\qquad
   \Pi=e^{-\tau}Q .
\]
If \(I=[\tau_1,\tau_2]\), then \(e^{-\tau}\) is bounded above and below by
positive constants on \(I\).  Therefore
\[
   \int_I\int_{B_R}|V|^2\,dy\,d\tau
   \le C(I)\int_{-e^{-\tau_1}}^{-e^{-\tau_2}}
       \int_{B_{C(I)R}}|U|^2\,dx\,dt,
\]
and similarly
\[
   \int_I\int_{B_R}|\nabla_yV|^2\,dy\,d\tau
   \le C(I)\int_{-e^{-\tau_1}}^{-e^{-\tau_2}}
       \int_{B_{C(I)R}}|\nabla_xU|^2\,dx\,dt.
\]
For the pressure, subtract a spatial-average gauge \(a(t)\) for \(Q\) on the
corresponding \(x\)-ball and set \(b(\tau)=e^{-\tau}a(-e^{-\tau})\).  Then
\(\Pi-b=e^{-\tau}(Q-a)\), and the same bounded-weight change of variables
gives \(\Pi-b\in L^{3/2}(B_R\times I)\).
\end{proof}

\begin{remark}\label{rem:no-global-tightness}
\Cref{lem:centered-local-energy} is local.  It does not give uniform
\(L^3\)-tightness in \(y\), nor a global \(L^\infty_\tau L^3_y\) bound.  This
is why the residual class is needed.
\end{remark}

\subsection{Normalized Seregin ancient limits}

\begin{definition}[Normalized Seregin ancient limit]\label{def:seregin-limit}
A smooth bounded solution \((V,\Pi)\) of \eqref{eq:centered} is a normalized
Seregin ancient limit generated by a Paper--0-admissible Type I source
\((u,p,z_*)\) if it arises from the blow-up construction of
\cref{thm:extraction} along an admissible Paper--0 concentration sequence and
satisfies the velocity nonvanishing condition \eqref{eq:nonvanishing}.
\end{definition}

\begin{definition}[Seregin collection]\label{def:seregin-collection}
A Seregin collection \(\calS\) is any nonempty collection of normalized
Seregin ancient limits generated by Paper--0-admissible Type I sources.  When
needed, \(\calS(u,p,z_*)\) denotes a collection associated with a fixed source.
\end{definition}

\begin{remark}
The nonvanishing condition is part of the normalization, not an additional
hypothesis.  It is inherited from the selected Paper--0 velocity
concentration and the no-escape condition through
\cref{lem:scale-selection,prop:nonzero}.  Thus every normalized Seregin
ancient limit is nonzero in the velocity sense of \eqref{eq:nonvanishing}.
\end{remark}

\begin{proof}[Proof of \cref{thm:extraction}]
Choose an admissible Paper--0 concentration sequence
\(\lambda_k\downarrow0\) from \cref{def:source}, and form the rescaled
sequence \eqref{eq:rescaled-sequence}.  The compact ancient limit is
\cref{prop:ancient-limit}; the strong and weak convergence are
\cref{prop:compactness}.  The no-escape argument in
\cref{lem:scale-selection}, followed by strong \(L^3\)-compactness away from
\(t=0\), gives nontriviality through \cref{prop:nonzero}.  The pointwise
Type I bound is \cref{cor:classical-type-I-bound}, and the centered equation,
smoothness, boundedness, and velocity nonvanishing are
\cref{lem:renormalized-equation,lem:smoothness,lem:renorm-bounds}.
\end{proof}

\section{Admissible classes and the residual class}\label{sec:classes}

\begin{definition}[Admissible branch parameter sets]
\label{def:admissible-parameters}
An admissible swirl parameter set is a nonempty set
\(I\subset(0,\infty)\) of exponents \(\gamma\) for which the corresponding
controlled-swirl Liouville criterion is assumed or proved.  An admissible decay
parameter set is a nonempty set \(J\subset(0,\infty)\) of exponents \(m\) for
which the corresponding weighted-decay Liouville criterion is assumed or
proved.
\end{definition}

Fix a Seregin collection \(\calS\) and admissible parameter sets
\[
   I,J\subset(0,\infty).
\]
All classes are subsets of \(\calS\), and complements are taken relative to
\(\calS\).  Structural conditions involving the unrenormalized ancient field
refer to
\[
   U(x,t)=(-t)^{-1/2}V(x/\sqrt{-t},-\log(-t)).
\]

\begin{definition}[Small-amplitude class]\label{def:small-branch}
Let \(\varepsilon_0>0\) be the small-data Liouville constant for bounded
ancient solutions of \eqref{eq:centered}.  We say
\(V\in\calX_{\mathrm{small}}\) if
\[
   \|V\|_{L^\infty(\R^3\times\R)}\le \varepsilon_0.
\]
\end{definition}

\begin{definition}[Stationary \(L^3\) class]\label{def:stationary-branch}
We say \(V\in\calX_{\mathrm{stat}\text{-}L^3}\) if
\[
   V(y,\tau)=W(y)
   \quad\hbox{and}\quad
   W\in L^3(\R^3).
\]
\end{definition}

\begin{definition}[Uniformly \(L^3\)-tight class]\label{def:tight-branch}
We say \(V\in\calX_{L^3\mathrm{-tight}}\) if
\begin{equation}\label{eq:L3-bound}
   \sup_{\tau\in\R}\|V(\tau)\|_{L^3(\R^3)}<\infty
\end{equation}
and
\begin{equation}\label{eq:L3-tightness}
   \lim_{R\to\infty}
   \sup_{\tau\in\R}\int_{|y|>R}|V(y,\tau)|^3\,dy=0 .
\end{equation}
\end{definition}

\begin{definition}[Structured/decay class]\label{def:structured-decay-branch}
A normalized Seregin limit \(V\) belongs to \(\calX_{\mathrm{sd}}(I,J)\) if
its physical ancient representative \(U\) satisfies at least one of the
following precisely stated conditions:
\begin{enumerate}[label=(\roman*)]
\item \(U\) is axisymmetric without swirl;
\item \(U\) is axisymmetric and, for some \(\gamma\in I\), its centered
representative satisfies the controlled-swirl bound
\[
   \sup_{\rho>0,\ Z\in\R,\ \tau\in\R}
   \rho^\gamma |V_\theta(\rho,Z,\tau)|<\infty .
\]
Equivalently, in physical variables,
\[
   \sup_{r>0,\ z\in\R,\ t<0}
   (-t)^{(1-\gamma)/2} r^\gamma |U_\theta(r,z,t)|<\infty .
\]
\item \(U\) is axisymmetric and
\[
   rU_\theta\in L_t^\infty L_x^p(\R^3\times(-\infty,0))
\]
for some \(1\le p<\infty\);
\item \(U\) is axisymmetric, periodic in the axial variable, and
\[
   rU_\theta\in L^\infty(\R^3\times(-\infty,0));
\]
\item \(V\) satisfies the weighted decay condition
\begin{equation}\label{eq:weighted-decay}
   \sup_{\tau\in\R}
   \int_{\R^3}(1+|y|^2)^m|V(y,\tau)|^2\,dy<\infty
\end{equation}
for some \(m\in J\).
\end{enumerate}
\end{definition}

\begin{lemma}[Structural quantities under centering]
\label{lem:structure-centering}
The centered change of variables preserves axisymmetry and absence of swirl.
If \((r,z,t)\) and \((\rho,Z,\tau)\) are related by
\[
   \rho=\frac r{\sqrt{-t}},\qquad Z=\frac z{\sqrt{-t}},
   \qquad \tau=-\log(-t),
\]
then the swirl variables satisfy
\[
   \rho V_\theta(\rho,Z,\tau)=rU_\theta(r,z,t).
\]
The controlled-swirl equivalence in
\cref{def:structured-decay-branch} follows from
\[
   \rho^\gamma |V_\theta(\rho,Z,\tau)|
   =
   (-t)^{(1-\gamma)/2}r^\gamma |U_\theta(r,z,t)|.
\]
\end{lemma}

\begin{proof}
The map \(x\mapsto y=x/\sqrt{-t}\) is radial in the cylindrical variables and
does not mix the \(e_r,e_\theta,e_z\) frame.  Since
\(V(y,\tau)=\sqrt{-t}\,U(x,t)\), the \(\theta\)-components obey
\(V_\theta=\sqrt{-t}\,U_\theta\), and the displayed identities follow.
\end{proof}

\begin{definition}[Residual Seregin class]\label{def:residual}
The residual class is
\[
   \calR(\calS;I,J)
   =
   \calS
   \setminus
   \left(
      \calX_{\mathrm{small}}
      \cup
      \calX_{\mathrm{stat}\text{-}L^3}
      \cup
      \calX_{L^3\mathrm{-tight}}
      \cup
      \calX_{\mathrm{sd}}(I,J)
   \right).
\]
\end{definition}

\begin{remark}
The non-residual classes are not intended to be disjoint.  The residual class
is the complement of their union, so overlaps among closed branches cause no
ambiguity.
\end{remark}

\begin{proposition}[Exhaustion by admissible classes]\label{prop:exhaustion}
Every \(V\in\calS\) belongs either to one of the four non-residual classes or
to \(\calR(\calS;I,J)\).
\end{proposition}

\begin{proof}
Let
\[
   \calU=
      \calX_{\mathrm{small}}
      \cup
      \calX_{\mathrm{stat}\text{-}L^3}
      \cup
      \calX_{L^3\mathrm{-tight}}
      \cup
      \calX_{\mathrm{sd}}(I,J).
\]
For \(V\in\calS\), either \(V\in\calU\), in which case it lies in at least one
non-residual class, or \(V\in\calS\setminus\calU=\calR(\calS;I,J)\).
\end{proof}

\section{Liouville hypotheses and non-residual exclusions}\label{sec:liouville}

\begin{lemma}[Centered Stokes kernel estimates]
\label{lem:centered-stokes-estimates}
Let
\[
   L=\Delta-\frac12y\cdot\nabla-\frac12,
   \qquad S(s)=e^{sL}.
\]
There is a universal constant \(C_0\) such that, for all \(s>0\),
\[
   \|S(s)f\|_{L^\infty}\le e^{-s/2}\|f\|_{L^\infty}
\]
and
\[
   \|S(s)\mathbb P\nabla\cdot F\|_{L^\infty}
   \le C_0 e^{-s/2}(1-e^{-s})^{-1/2}\|F\|_{L^\infty}.
\]
\end{lemma}

\begin{proof}
The first estimate follows from the explicit Ornstein--Uhlenbeck kernel:
\[
   S(s)f(y)=e^{-s/2}
   \int_{\R^3}G_{1-e^{-s}}(e^{-s/2}y-z)f(z)\,dz,
\]
where \(G_a\) is the heat kernel with variance \(a\).  The kernel has total
mass one.  For the projected-divergence estimate, write the corresponding
kernel as the Oseen projected-gradient kernel after the same change of
variables.  More explicitly, the semigroup for
\(\Delta-\frac12y\cdot\nabla\) is
\[
   (e^{s(\Delta-\frac12y\cdot\nabla)}g)(y)
   =
   (e^{(1-e^{-s})\Delta}g)(e^{-s/2}y).
\]
Applying this to the distribution \(g=\mathbb P\nabla\cdot F\) and using the
projected-gradient Oseen kernel estimate
\[
   \|e^{\theta\Delta}\mathbb P\nabla\cdot F\|_\infty
   \le C\theta^{-1/2}\|F\|_\infty,\qquad \theta>0,
\]
gives the displayed estimate with \(\theta=1-e^{-s}\), together with the
extra factor \(e^{-s/2}\) from the zeroth-order term in \(L\).  The
projected-gradient Oseen estimate follows from the \(L^1\)-bound
\[
   \|\nabla O(\cdot,\theta)\|_{L^1}\le C\theta^{-1/2}
\]
for the gradient of the Oseen kernel \(O(\cdot,\theta)\).  This is the
bounded-data Stokes estimate used by Kato \cite{Kato1984}; it does not require
boundedness of the Leray projector itself on \(L^\infty\).
\end{proof}

\begin{lemma}[Mild representation for normalized limits]
\label{lem:mild-representation}
Every normalized Seregin ancient limit satisfies, for all \(\Theta>0\),
\[
   V(\tau)
   =
   S(\Theta)V(\tau-\Theta)
   -
   \int_0^\Theta
   S(\Theta-\sigma)\mathbb P\nabla\cdot
   (V\otimes V)(\tau-\Theta+\sigma)\,d\sigma .
\]
\end{lemma}

\begin{proof}
Normalized Seregin limits are smooth and bounded by
\cref{lem:smoothness,lem:renorm-bounds}.  On a finite interval
\([\tau-\Theta,\tau]\), the tensor \(V\otimes V\) is bounded, so
\cref{lem:centered-stokes-estimates} makes the Duhamel integral finite in
\(L^\infty\).  In the sense of solenoidal distributions, applying the Leray
projector to \eqref{eq:centered} removes the pressure and gives
\[
   \partial_\tau V=LV-\mathbb P\nabla\cdot(V\otimes V)
\]
on every finite interval.  Since \(V\) is smooth, the map
\[
   s\mapsto S(\tau-s)V(s)
\]
is differentiable as a tempered distribution for \(\tau-\Theta<s<\tau\), and
\[
   \frac{d}{ds}\bigl(S(\tau-s)V(s)\bigr)
   =
   -S(\tau-s)\mathbb P\nabla\cdot(V\otimes V)(s).
\]
Integrating this identity from \(\tau-\Theta\) to \(\tau\) gives the stated
Duhamel formula.
\end{proof}

\begin{lemma}[Small bounded ancient Liouville theorem]\label{lem:small-liouville}
There is \(\varepsilon_0>0\) such that every bounded mild ancient solution of
\eqref{eq:centered} with
\[
   \|V\|_{L^\infty(\R^3\times\R)}\le\varepsilon_0
\]
is identically zero.
\end{lemma}

\begin{proof}
Write the projected centered equation as
\[
   \partial_\tau V=LV-\mathbb P\nabla\cdot(V\otimes V),
   \qquad
   L=\Delta-\frac12 y\cdot\nabla-\frac12.
\]
Let \(S(s)=e^{sL}\).  By \cref{lem:centered-stokes-estimates}, the nonlinear
kernel is integrable on \((0,\infty)\); set
\[
   A=C_0\int_0^\infty e^{-s/2}(1-e^{-s})^{-1/2}\,ds<\infty .
\]
For fixed \(\tau\) and \(T>0\), the mild identity on \([\tau-T,\tau]\) gives
\[
   V(\tau)
   =
   S(T)V(\tau-T)
   -
   \int_0^T
   S(T-\sigma)\mathbb P\nabla\cdot(V\otimes V)(\tau-T+\sigma)\,d\sigma .
\]
If \(M=\|V\|_{L^\infty(\R^3\times\R)}\), then
\[
   \|V(\tau)\|_\infty\le e^{-T/2}M+AM^2.
\]
Letting \(T\to\infty\) and taking the supremum in \(\tau\) gives
\[
   M\le AM^2.
\]
If \(0<M<1/A\), this is impossible.  Choose \(\varepsilon_0<1/A\); then
\(M\le\varepsilon_0\) implies \(M=0\).
\end{proof}

\begin{lemma}[Stationary centered profiles]\label{lem:stationary-profile}
If \(V(y,\tau)=W(y)\) is a smooth solution of \eqref{eq:centered}, then after
subtracting a function of \(\tau\) from the pressure, \(W\) solves
\[
   (W\cdot\nabla)W+\nabla P
   =
   \Delta W-\frac12(W+y\cdot\nabla W),
   \qquad \divergence W=0.
\]
\end{lemma}

\begin{proof}
Since \(V\) is independent of \(\tau\), \(\partial_\tau V=0\).  Set
\[
   F(y)=
   \Delta W-\frac12(W+y\cdot\nabla W)-(W\cdot\nabla)W .
\]
The centered equation gives \(\nabla_y\Pi(\cdot,\tau)=F\) for a.e. \(\tau\).
Fix one such time \(\tau_0\) and put \(P(y)=\Pi(y,\tau_0)\).  Then
\(\nabla_y(\Pi(\cdot,\tau)-P)=0\) for a.e. \(\tau\), so
\(\Pi(y,\tau)-P(y)=b(\tau)\) is spatially constant.  Subtracting this gauge
from \(\Pi\) leaves the equation unchanged and yields the displayed
stationary self-similar equation for \(W\).
\end{proof}

\begin{theorem}[Small and stationary branch exclusions]
\label{thm:classical-closed-branches}
Let \(V\) be a normalized Seregin ancient limit.  Then
\[
   V\notin \calX_{\mathrm{small}},
   \qquad
   V\notin \calX_{\mathrm{stat}\text{-}L^3}.
\]
\end{theorem}

\begin{proof}
If \(V\in\calX_{\mathrm{small}}\), then
\cref{lem:mild-representation,lem:small-liouville} give \(V\equiv0\),
contradicting the velocity nonvanishing condition \eqref{eq:nonvanishing}.

If \(V\in\calX_{\mathrm{stat}\text{-}L^3}\), then
\(V(y,\tau)=W(y)\) with \(W\in L^3(\R^3)\).  By
\cref{lem:stationary-profile}, \(W\) satisfies the stationary self-similar
profile equation after a pressure gauge change.
The theorem of Ne\v{c}as--R\r{u}\v{z}i\v{c}ka--\v{S}ver\'ak
\cite{NRS1996} gives \(W\equiv0\).  Again
\eqref{eq:nonvanishing} gives the contradiction.
\end{proof}

\begin{hypothesis}[Tight-class Liouville hypothesis]\label{hyp:tight-liouville}
No nonzero normalized Seregin ancient limit belongs to
\(\calX_{L^3\mathrm{-tight}}\).
\end{hypothesis}

\begin{remark}
This hypothesis is logically independent of the blow-up extraction.  It is the
form one obtains from a minimal-element reduction together with compact-hull
equilibrium rigidity in the uniformly \(L^3\)-tight class.
\end{remark}

\begin{hypothesis}[Structured/decay Liouville hypothesis]
\label{hyp:structured-decay-liouville}
For the admissible parameter sets \(I,J\), no nonzero normalized Seregin
ancient limit belongs to \(\calX_{\mathrm{sd}}(I,J)\).
\end{hypothesis}

\begin{remark}
The structured/decay hypothesis includes the axisymmetric no-swirl theorem, the
\(L_t^\infty L_x^p\)-swirl theorem, the periodic bounded-swirl theorem, and
the additional weighted-decay or controlled-swirl Liouville criteria supplied
by a structured-branch analysis.  The published axisymmetric components used
as model cases are \cite{KNSS2009,LZZ2017,LRZ2022}.
\end{remark}

\begin{proposition}[Non-residual class exclusion]
\label{prop:closed-branch-exclusion}
Assume \cref{hyp:tight-liouville,hyp:structured-decay-liouville}.  Then no
normalized Seregin ancient limit can belong to
\[
   \calX_{\mathrm{small}}
   \cup
   \calX_{\mathrm{stat}\text{-}L^3}
   \cup
   \calX_{L^3\mathrm{-tight}}
   \cup
   \calX_{\mathrm{sd}}(I,J).
\]
\end{proposition}

\begin{proof}
The small and stationary \(L^3\) classes are excluded by
\cref{thm:classical-closed-branches}.  The tight class is excluded by
\cref{hyp:tight-liouville}.  The structured/decay class is excluded by
\cref{hyp:structured-decay-liouville}.
\end{proof}

\section{The residual-class criterion}\label{sec:criterion}

\begin{hypothesis}[No residual Seregin class]\label{hyp:no-residual}
For every Seregin collection \(\calS\) generated by Paper--0-admissible
finite-energy Type I sources,
\[
   \calR(\calS;I,J)=\varnothing .
\]
\end{hypothesis}

\begin{hypothesis}[Residual rigidity]\label{hyp:residual-rigidity}
For every Seregin collection \(\calS\) generated by Paper--0-admissible
finite-energy Type I sources, every element of \(\calR(\calS;I,J)\) is
identically zero.
\end{hypothesis}

\begin{proposition}[Equivalence on normalized Seregin limits]
\label{prop:residual-equivalence}
On normalized Seregin collections, \cref{hyp:no-residual} and
\cref{hyp:residual-rigidity} are equivalent.
\end{proposition}

\begin{proof}
If \cref{hyp:no-residual} holds, residual rigidity is vacuous.  Conversely,
assume \cref{hyp:residual-rigidity}.  If some normalized collection had
\(V\in\calR(\calS;I,J)\), residual rigidity would give \(V\equiv0\), directly
contradicting the positive velocity normalization \eqref{eq:nonvanishing}.
Hence the residual class is empty.
\end{proof}

\begin{proof}[Proof of \cref{thm:residual-criterion}]
Suppose that \((u,p,z_*)\) is a Paper--0-admissible finite-energy Type I
source.  By \cref{thm:extraction}, there exists a normalized nonzero Seregin
ancient limit \(V\).  Let \(\calS\) be a Seregin collection containing \(V\).

By \cref{prop:exhaustion}, either \(V\) belongs to a non-residual class or
\(V\in\calR(\calS;I,J)\).  The first alternative is impossible by
\cref{prop:closed-branch-exclusion}.  The second alternative is impossible by
\cref{hyp:no-residual}.  Hence the assumed Paper--0-admissible Type I source
cannot exist.
\end{proof}

\begin{corollary}[Replacement residual rigidity criterion]
Assume \cref{hyp:tight-liouville,hyp:structured-decay-liouville,hyp:residual-rigidity}.
Then no Paper--0-admissible finite-energy Type I source exists.
\end{corollary}

\begin{proof}
Use \cref{prop:residual-equivalence} to replace
\cref{hyp:residual-rigidity} by \cref{hyp:no-residual}, then apply
\cref{thm:residual-criterion}.
\end{proof}

The preceding theorem isolates the obstruction to Paper--0-admissible
finite-energy Type I exclusion as a residual statement for normalized Seregin
ancient limits.  The compactness construction, pressure normalization, and
branch exhaustion are contained in the present paper, with the local
velocity concentration input stated explicitly in
\cref{prop:paper0-velocity-concentration}.  The remaining analytic hypotheses are
precisely the tight-class Liouville theorem, the structured/decay Liouville
theorem, and the no-residual assertion.

\appendix

\section{Absence of Type II cascade symmetries}\label{app:no-cascade}

\begin{lemma}[Residual logarithmic-time translation]\label{lem:tau-translation}
If \(V\) solves \eqref{eq:centered}, then
\[
   V_{\tau_0}(y,\tau)=V(y,\tau+\tau_0)
\]
also solves \eqref{eq:centered}.  If
\[
   U^\lambda(x,t)=\lambda U(\lambda x,\lambda^2t),
\]
then the renormalized field associated with \(U^\lambda\) is
\[
   V^\lambda(y,\tau)=V(y,\tau-2\log\lambda).
\]
\end{lemma}

\begin{proof}
The first statement follows because \eqref{eq:centered} is autonomous in
\(\tau\).  For the second, put \(s=-t=e^{-\tau}\).  Then
\[
   \sqrt{s}\,U^\lambda(y\sqrt{s},-s)
   =
   \lambda\sqrt{s}\,U(\lambda y\sqrt{s},-\lambda^2s)
   =
   V(y,-\log(\lambda^2s))
   =
   V(y,\tau-2\log\lambda).
\]
\end{proof}

\begin{proposition}[Failure of generic dilation and translation symmetries]
\label{prop:no-cascade}
The centered equation \eqref{eq:centered} is autonomous in \(\tau\), but it is
not invariant under
\[
   V(y,\tau)\mapsto \alpha V(\alpha y,\alpha^2\tau)
\]
unless \(\alpha=1\).  It is also not invariant under constant translations
\[
   V(y,\tau)\mapsto V(y-a,\tau)
\]
unless \(a=0\).
\end{proposition}

\begin{proof}
For scaling, set \(Y=\alpha y\), \(T=\alpha^2\tau\), and
\(\widetilde V(y,\tau)=\alpha V(Y,T)\),
\(\widetilde\Pi(y,\tau)=\alpha^2\Pi(Y,T)\).  The non-drift part scales with
factor \(\alpha^3\), while
\[
   -\frac12(\widetilde V+y\cdot\nabla\widetilde V)
   =
   -\frac{\alpha}{2}(V+Y\cdot\nabla_YV)(Y,T).
\]
Using the equation for \(V\), the residual for \(\widetilde V\) is
\[
   \frac{\alpha-\alpha^3}{2}(V+Y\cdot\nabla_YV)(Y,T).
\]
If \(\alpha\ne1\), this forces \(V+Y\cdot\nabla V=0\).  Along each ray,
\[
   \frac{d}{dr}\bigl(rV(r\omega,T)\bigr)=0.
\]
Boundedness and smoothness at the origin force the constant to be zero, hence
\(V\equiv0\).  Thus a nonzero centered profile has no such scaling symmetry.

For translations, all translation-invariant terms transform by composition,
but the drift satisfies
\[
   y\cdot\nabla V(y-a,\tau)
   =
   (y-a)\cdot\nabla V(y-a,\tau)+a\cdot\nabla V(y-a,\tau).
\]
The extra term is absent from the original equation.  Thus spatial
translation is not a symmetry of the fixed centered equation unless \(a=0\),
apart from exceptional fixed solutions satisfying \(a\cdot\nabla V\equiv0\).
\end{proof}

\begin{thebibliography}{99}

\bibitem{AlbrittonBarker2019}
D. Albritton and T. Barker,
\emph{On local Type I singularities of the Navier--Stokes equations and
Liouville theorems},
J. Math. Fluid Mech. \textbf{21} (2019), article no. 43,
doi:\url{https://doi.org/10.1007/s00021-019-0448-z}.

\bibitem{Aubin1963}
J.-P. Aubin,
\emph{Un theoreme de compacite},
C. R. Acad. Sci. Paris \textbf{256} (1963), 5042--5044.

\bibitem{CKN1982}
L. Caffarelli, R. Kohn, and L. Nirenberg,
\emph{Partial regularity of suitable weak solutions of the Navier--Stokes equations},
Comm. Pure Appl. Math. \textbf{35} (1982), no. 6, 771--831,
doi:\url{https://doi.org/10.1002/cpa.3160350604}.

\bibitem{ESS2003}
L. Escauriaza, G. Seregin, and V. \v{S}verak,
\emph{\(L_{3,\infty}\)-solutions of Navier--Stokes equations and backward uniqueness},
Russian Math. Surveys \textbf{58} (2003), no. 2, 211--250,
doi:\url{https://doi.org/10.1070/RM2003v058n02ABEH000609}.

\bibitem{Kato1984}
T. Kato,
\emph{Strong \(L^p\)-solutions of the Navier--Stokes equation in \(\mathbb R^m\), with applications to weak solutions},
Math. Z. \textbf{187} (1984), no. 4, 471--480,
doi:\url{https://doi.org/10.1007/BF01174182}.

\bibitem{KNSS2009}
G. Koch, N. Nadirashvili, G. Seregin, and V. \v{S}verak,
\emph{Liouville theorems for the Navier--Stokes equations and applications},
Acta Math. \textbf{203} (2009), no. 1, 83--105.

\bibitem{Lin1998}
F.-H. Lin,
\emph{A new proof of the Caffarelli--Kohn--Nirenberg theorem},
Comm. Pure Appl. Math. \textbf{51} (1998), no. 3, 241--257,
doi:\url{https://doi.org/10.1002/(SICI)1097-0312(199803)51:3<241::AID-CPA2>3.0.CO;2-A}.

\bibitem{Leray1934}
J. Leray,
\emph{Sur le mouvement d'un liquide visqueux emplissant l'espace},
Acta Math. \textbf{63} (1934), 193--248.

\bibitem{LionsMagenes}
J.-L. Lions and E. Magenes,
\emph{Non-homogeneous boundary value problems and applications. Vol. I},
Springer-Verlag, 1972.

\bibitem{LRZ2022}
Z. Lei, X. Ren, and Q. S. Zhang,
\emph{A Liouville theorem for the axially symmetric Navier--Stokes equations},
J. Funct. Anal. \textbf{283} (2022), no. 3, 109522.

\bibitem{LZZ2017}
Z. Lei, Q. Zhang, and N. Zhao,
\emph{Improved Liouville theorems for axially symmetric Navier--Stokes equations},
Sci. Sin. Math. \textbf{47} (2017), no. 10, 1183--1198.

\bibitem{NRS1996}
J. Necas, M. Ruzicka, and V. \v{S}verak,
\emph{On Leray's self-similar solutions of the Navier--Stokes equations},
Acta Math. \textbf{176} (1996), no. 2, 283--294.

\bibitem{Paper0LocalConcentration}
G. Duran Ballester,
\emph{Local CKN concentration and the Type I/Type II dichotomy for the
Navier--Stokes equations},
preprint, 2026.

\bibitem{Seregin2012}
G. Seregin,
\emph{A certain necessary condition of potential blow up for Navier--Stokes equations},
Comm. Math. Phys. \textbf{312} (2012), no. 3, 833--845,
doi:\url{https://doi.org/10.1007/s00220-011-1391-x}.

\bibitem{Serrin1962}
J. Serrin,
\emph{On the interior regularity of weak solutions of the Navier--Stokes equations},
Arch. Rational Mech. Anal. \textbf{9} (1962), 187--195,
doi:\url{https://doi.org/10.1007/BF00253344}.

\bibitem{Simon1987}
J. Simon,
\emph{Compact sets in the space \(L^p(0,T;B)\)},
Ann. Mat. Pura Appl. (4) \textbf{146} (1987), 65--96,
doi:\url{https://doi.org/10.1007/BF01762360}.

\bibitem{Sohr2001}
H. Sohr,
\emph{The Navier--Stokes equations. An elementary functional analytic approach},
Birkhauser, 2001.

\bibitem{Stein1970}
E. M. Stein,
\emph{Singular integrals and differentiability properties of functions},
Princeton University Press, 1970.

\end{thebibliography}

\end{document}
